\section{From a counterexample to a strict subset inequality}
\label{sec:reduction}

We first convert a hypothetical counterexample into a graph whose first
image is larger than the average of a set and its second image.  Throughout
this section, $F=F_D$ denotes the ordinary union of the out-neighborhoods
in a specified base graph $D$; in particular, $F(S)$ need not be disjoint
from $S$.

\begin{proposition}[Counterexample reduction]
\label{prop:counterexample-reduction}
If Theorem~\ref{thm:main} has a counterexample, then there is a finite
nonempty oriented graph $G$ in which every vertex has positive indegree
and
\begin{equation}
\label{eq:strict-subset-growth}
  \abs{U}+\abs{F_G^2(U)}<2\abs{F_G(U)}
  \qquad\text{for every }\varnothing\ne U\subsetneq V(G).
\end{equation}
\end{proposition}

There are two steps.  Minimality first gives a strict deficit for subsets
of a counterexample.  Replacing each vertex by a sufficiently large
transitive tournament then turns that deficit into
\eqref{eq:strict-subset-growth}.

The arc-deletion argument below adapts the proof idea of
\cite[Lemma~4]{Seacrest2019}.  Specializing Seacrest's parameter $\lambda$
to $1$, his statement gives
$\abs{F^2(S)\setminus F(S)}<\abs{F(S)}$ in an edge-minimal
counterexample when $F(S)$ is nonempty.  Lemma~\ref{lem:subset-deficit}
replaces this right-hand bound by $\abs{F(S)\setminus S}$ for nonempty
proper $S$, under the stronger choice of minimum order and then minimum
arc count.  We prove this strengthened form in full.

\begin{lemma}[Subset deficit]
\label{lem:subset-deficit}
Let $D$ be a finite nonempty oriented graph such that
$\abs{N_1^+(v)}>\abs{N_2^+(v)}$ for every $v\in V(D)$.
Choose $D$ with minimum order and, subject to that, minimum number of
arcs among all graphs with this property.  Then $D$ is strongly connected,
every vertex has positive indegree, and
\begin{equation}
\label{eq:subset-deficit}
  \abs{F^2(S)\setminus F(S)}<\abs{F(S)\setminus S}
  \qquad\text{for every }\varnothing\ne S\subsetneq V(D).
\end{equation}
\end{lemma}

\begin{proof}
A sink satisfies the second-neighborhood inequality, so $D$ has no sink
and has more than one vertex.  If $D$ were not strongly connected, take
a strongly connected component with no outgoing arc to another component.
All first and second out-neighborhoods of its vertices lie in that
component.  Its induced graph would therefore be a smaller counterexample,
contrary to the choice of $D$.  Thus $D$ is strongly connected, and every
vertex has an in-neighbor.

Fix a nonempty proper subset $S$ and put
\[
  I=F(S)\cap S,\qquad E=F(S)\setminus S,
  \qquad g(T)=\abs{F(T)\setminus(I\cup T)}\quad(T\subseteq E).
\]
Strong connectivity gives $E\ne\varnothing$.  We will enlarge $T$ from
$\varnothing$ to $E$ so that the increase in $g(T)$ is strictly smaller
than the number of newly adjoined vertices at every step.

For $T\subsetneq E$, let $D_T$ be obtained from the original graph $D$
by deleting all arcs from $S$ to $E\setminus T$.  Since every vertex of
$E\setminus T$ receives an arc from $S$, at least one arc is deleted.
Minimality of the number of arcs gives a vertex $v$ satisfying
$\abs{N_{1,D_T}^+(v)}\le\abs{N_{2,D_T}^+(v)}$.
If $v$ had lost no first out-neighbor, deleting arcs could only shrink
its exact second out-neighborhood, so it would still violate this
inequality.  Consequently $v\in S$, and the set
\[
  Q=F(\{v\})\cap(E\setminus T)
\]
is nonempty.  It is exactly the set of first out-neighbors that $v$ loses.
The original first out-neighborhood of $v$ is contained in $I\cup T\cup Q$,
and the remaining first out-neighbors lie in $I\cup T$.

We compare the exact second out-neighborhoods before and after deletion.
A two-step path from $v$ in $D_T$ through a vertex of $I\subseteq S$
ends in $I\cup T$, since all arcs from $S$ to $E\setminus T$ were
deleted.  A path through a vertex of $T$ ends in $F(T)$.  Thus
\begin{equation}
\label{eq:deleted-path-endpoints}
  F_{D_T}^2(\{v\})\subseteq I\cup T\cup F(T).
\end{equation}
Every newly counted exact second out-neighbor must have been a first
out-neighbor in $D$: deletion creates no paths, and $v$ itself is never
an exact second out-neighbor.  Such a vertex lies in $Q$, and
\eqref{eq:deleted-path-endpoints}, together with
$Q\cap(I\cup T)=\varnothing$, places it in $Q\cap F(T)$.

On the other hand, define
\[
  C=F(Q)\setminus\bigl(I\cup T\cup Q\cup F(T)\bigr).
\]
Every vertex of $C$ was reached from $v$ in two steps through $Q$, and
none was a first out-neighbor.  Also $v\notin F(Q)$: an arc from a
vertex of $Q$ back to $v$ would oppose the arc from $v$ to that vertex.
Hence $C\subseteq N_{2,D}^+(v)$.  By
\eqref{eq:deleted-path-endpoints}, no vertex of $C$ remains reachable
in two steps in $D_T$.

Write $d_1=\abs{N_{1,D}^+(v)}$ and
$d_2=\abs{N_{2,D}^+(v)}$.  The gains and losses just identified imply
\[
  d_1-\abs{Q}
  \le \abs{N_{2,D_T}^+(v)}
  \le d_2-\abs{C}+\abs{Q\cap F(T)}.
\]
The change in $g$ is exact: adjoining $Q$ removes the points
$Q\cap F(T)$ from the previously counted image and adds precisely $C$.
Since $d_2<d_1$, it follows that
\begin{equation}
\label{eq:deficit-increment}
  g(T\cup Q)-g(T)
  =\abs{C}-\abs{Q\cap F(T)}
  \le \abs{Q}+d_2-d_1
  <\abs{Q}.
\end{equation}

Starting with $T=\varnothing$, repeat this step until $T=E$.
Each deletion is performed afresh on the original $D$, so the same
minimality argument applies at every step.  Each $Q$ is a nonempty
subset of the part of $E$ not yet adjoined; hence the process terminates.
Summing \eqref{eq:deficit-increment} and using $g(\varnothing)=0$
gives $g(E)<\abs{E}$.  Finally,
$F(S)=I\cup E$ and $F(I)\subseteq F(S)$, so
\[
  F^2(S)\setminus F(S)
  =\bigl(F(I)\cup F(E)\bigr)\setminus(I\cup E)
  =F(E)\setminus(I\cup E).
\]
Its cardinality is $g(E)$, proving \eqref{eq:subset-deficit}.
\end{proof}

We next amplify the integer deficit in \eqref{eq:subset-deficit}.
Each unit of the base deficit contributes the entire block size; the
remaining internal contribution costs at most one unit per base vertex.
The construction below is the lexicographic product $D[T_m]$, where $T_m$
is a transitive tournament.  The same graph-product construction appears
in \cite[Theorem~4.3]{BrantnerBrockmanKaySnively2009}; here we need an
estimate for every nonempty proper subset, rather than only for individual
vertices.

\begin{lemma}[Transitive-tournament blowup]
\label{lem:blowup}
Let $D$ be an oriented graph on $n\ge1$ vertices, all of positive
indegree, satisfying \eqref{eq:subset-deficit} for every nonempty proper
subset of $V(D)$.  For an integer $m>n$, replace each vertex by a
transitive tournament on $m$ vertices.  For every arc $v\to w$ of $D$,
put all arcs from the block at $v$ to the block at $w$, and put no other
arcs between distinct blocks.  The resulting graph $G$ is oriented,
all its indegrees are positive, and it satisfies
\eqref{eq:strict-subset-growth}.
\end{lemma}

\begin{proof}
Label the vertices within each block by $1,\ldots,m$, with arcs from
smaller to larger labels.  Let $f$ be the image operator within this
transitive tournament.  For a nonempty set $T$ of labels with least
element $a$, we have $f(T)=\{a+1,\ldots,m\}$.  Thus
\begin{equation}
\label{eq:internal-block-bound}
  \abs{T}-\abs{f(T)}\le1,\qquad
  f^2(T)\subseteq f(T),\qquad
  \abs{T}-2\abs{f(T)}+\abs{f^2(T)}\le1.
\end{equation}
The construction produces an oriented graph.  Every vertex receives
arcs from an incoming base block, so all indegrees in $G$ are positive.
Call an arc external if it joins distinct blocks, and internal otherwise.

Fix $\varnothing\ne U\subsetneq V(G)$.  Write $T_v$ for the labels
selected by $U$ in the block at $v$, and define its support and a set
of base vertices by
\[
  S=\{v\in V(D):T_v\ne\varnothing\},\qquad
  H_0=F(S)\cup F^2(S).
\]
Every block indexed by $F(S)$ is filled by $F_G(U)$.  In each remaining
block the first image consists exactly of $f(T_v)$, so
\[
  \abs{F_G(U)}=m\abs{F(S)}
     +\sum_{v\in S\setminus F(S)}\abs{f(T_v)}.
\]
A two-step path using two external arcs ends in an $F^2(S)$-block.
One using exactly one external arc ends in an $F(S)$-block, regardless
of the order of its internal and external steps.  Paths using only
internal arcs contribute $f^2(T_v)$.  Therefore
\[
  \abs{F_G^2(U)}\le m\abs{H_0}
     +\sum_{v\in S\setminus H_0}\abs{f^2(T_v)}.
\]
Combining these estimates gives the sharper bound
\begin{align}
\label{eq:blowup-sharp}
  \abs{U}-2\abs{F_G(U)}+\abs{F_G^2(U)}
  &\le m\bigl(\abs{H_0}-2\abs{F(S)}\bigr)
       +\sum_{v\in S}\abs{T_v} \notag\\
  &\quad -2\sum_{v\in S\setminus F(S)}\abs{f(T_v)}
       +\sum_{v\in S\setminus H_0}\abs{f^2(T_v)}.
\end{align}
Since $F(S)\subseteq H_0$, extending the last sum to
$S\setminus F(S)$ only adds nonnegative terms.  This may count internal
second images in blocks already bounded by a full block in
\eqref{eq:blowup-sharp}; it remains a valid upper bound.  Now use
$\abs{T_v}\le m$ on $S\cap F(S)$ and
\eqref{eq:internal-block-bound} on $S\setminus F(S)$ to obtain
\begin{align*}
  \abs{U}-2\abs{F_G(U)}+\abs{F_G^2(U)}
  &\le m\bigl(\abs{H_0}-2\abs{F(S)}
                    +\abs{S\cap F(S)}\bigr)
       +\abs{S\setminus F(S)}\\
  &\le m\bigl(\abs{F^2(S)\setminus F(S)}
                    -\abs{F(S)\setminus S}\bigr)+n.
\end{align*}
If $S\subsetneq V(D)$, the parenthesized difference is an integer
strictly less than zero by \eqref{eq:subset-deficit}.  The last bound
is consequently at most $-m+n<0$.

It remains to consider $S=V(D)$, for which the subset-deficit hypothesis
does not apply.  Positive indegrees in $D$ give $F(S)=V(D)$, so
$F_G(U)=V(G)$.  Positive indegrees in $G$ then give
$F_G^2(U)=V(G)$ as well.  In this case
\[
  \abs{U}-2\abs{F_G(U)}+\abs{F_G^2(U)}
  =\abs{U}-\abs{V(G)}<0,
\]
because $U$ is proper.  This proves the required inequality for every
allowed $U$.
\end{proof}

\begin{proof}[Proof of Proposition~\ref{prop:counterexample-reduction}]
Choose a counterexample of minimum order and then minimum number of
arcs.  Lemma~\ref{lem:subset-deficit} supplies positive indegrees and
the subset deficit, so Lemma~\ref{lem:blowup}, with any integer
$m>\abs{V(D)}$, supplies the required graph $G$.
\end{proof}
